\documentclass[11pt]{article}
\usepackage{graphicx} % Required for inserting images
\usepackage[margin=2.5cm]{geometry}
\usepackage{listings}
\usepackage{hyperref}
\usepackage{array}
\usepackage{booktabs}
\usepackage{xcolor}
\usepackage{amsmath}
\usepackage{amssymb}
\usepackage{minted}
\usepackage{mdframed}
\usepackage{nameref}
\usepackage[spanish]{babel}
\definecolor{bgcolor}{rgb}{0.95, 0.95, 0.95}

\usepackage{xcolor}

\renewcommand{\figurename}{Figura}
\renewcommand{\contentsname}{Índice de Contenidos}
\renewcommand{\listingscaption}{Código}
\renewcommand{\tablename}{Tabla}

\definecolor{codegreen}{rgb}{0,0.6,0}
\definecolor{codegray}{rgb}{0.5,0.5,0.5}
\definecolor{codepurple}{rgb}{0.58,0,0.82}
\definecolor{backcolour}{rgb}{0.95,0.95,0.92}

\lstdefinestyle{mystyle}{
    backgroundcolor=\color{backcolour},   
    commentstyle=\color{codegreen},
    keywordstyle=\color{magenta},
    numberstyle=\tiny\color{codegray},
    stringstyle=\color{codepurple},
    basicstyle=\ttfamily\footnotesize,
    breakatwhitespace=false,         
    breaklines=true,                 
    captionpos=b,                    
    keepspaces=true,                 
    numbers=left,                    
    numbersep=5pt,                  
    showspaces=false,                
    showstringspaces=false,
    showtabs=false,                  
    tabsize=2,
    frame=single
}

\lstset{literate=
  {á}{{\'a}}1
  {é}{{\'e}}1
  {í}{{\'i}}1
  {ó}{{\'o}}1
  {ú}{{\'u}}1
  {Á}{{\'A}}1
  {É}{{\'E}}1
  {Í}{{\'I}}1
  {Ó}{{\'O}}1
  {Ú}{{\'U}}1
  {ñ}{{\~n}}1
  {ü}{{\"u}}1
  {Ü}{{\"U}}1
}

\lstset{style=mystyle}

\begin{document}

\begin{titlepage}
\begin{center}
    \begin{minipage}{0.4\textwidth}
        
\includegraphics[height=30px]{murcialog.png}
    \end{minipage}\hfill
    \begin{minipage}{0.55\textwidth}
        \hfill{
\includegraphics[height=40px]{layout_set_logo.png}}
    \end{minipage}

    \vspace{10px}

    \hrule

    \vspace{10px}

    {\Huge \textsc{Documentación \\ \vspace{10px} Proyecto AR}}

    \vspace{20px}
    
    Jiahao Zhu, NIE: X8650158L \\ Samuel Fernández, DNI: 49169424R

    \vspace{3px}
    
    Grupo (y subgrupo): 9.1.

    \vspace{15px}

    \today
\end{center}

\tableofcontents
\end{titlepage}

\newpage

\section{Introducción}

\section{Cuestiones Teóricas}

\subsection{Encaminamiento Intra-Dominio IPv4}

\subsubsection{Organización A}

\begin{enumerate}
    \item \textit{Muestre las tablas de rutas del Router A4 y comente los aspectos más relevantes. ¿Cuál es el camino óptimo para alcanzar la interfaz de R A2 que conecta con
la Organización B? ¿Por qué? ¿Cuántas alternativas hay
para alcanzarlo según la tabla de rutas? ¿Cuáles son y por qué?}

\vspace{7px}

\textsc{Contestación}: Tal y como refleja Figura \ref{fig:2.1.1.1}, para llegar a la subred de IP 172.84.18.224/30 desde el Router A4, el número mínimo de saltos es 2, y existen 3 vías distintas a seguir. Se trata de las siguientes rutas:

\begin{itemize}
    \item R A4 $\xrightarrow[\text{Se7/0}]{}$ 172.84.18.222 (R A6) $\xrightarrow[\text{Se3/0}]{}$ 172.84.18.209 (R A2)

    \item R A4 $\xrightarrow[\text{Se6/0}]{}$ 172.84.18.213 (R A1) $\xrightarrow[\text{Se3/0}]{}$ 172.84.18.194 (R A2)

    \item R A4 $\xrightarrow[\text{Se7/0}]{}$ 172.84.18.201 (R A3) $\xrightarrow[\text{Se6/0}]{}$ 172.84.18.205 (R A2)
\end{itemize}

\begin{figure}[h]
\begin{minted}[linenos, bgcolor=backcolour, frame=single, breaklines=true]{text}
Gateway of last resort is not set

     172.84.0.0/16 is variably subnetted, 12 subnets, 4 masks
R       172.84.16.0/23 [120/1] via 172.84.18.218, 00:00:12, Serial2/0
R       172.84.18.0/25 [120/1] via 172.84.18.201, 00:00:01, Serial3/0
R       172.84.18.128/26 [120/1] via 172.84.18.222, 00:00:26, Serial7/0
R       172.84.18.192/30 [120/1] via 172.84.18.201, 00:00:01, Serial3/0
R       172.84.18.196/30 [120/1] via 172.84.18.201, 00:00:01, Serial3/0
                         [120/1] via 172.84.18.213, 00:00:03, Serial6/0
C       172.84.18.200/30 is directly connected, Serial3/0
R       172.84.18.204/30 [120/1] via 172.84.18.213, 00:00:03, Serial6/0
R       172.84.18.208/30 [120/1] via 172.84.18.222, 00:00:26, Serial7/0
C       172.84.18.212/30 is directly connected, Serial6/0
C       172.84.18.216/30 is directly connected, Serial2/0
C       172.84.18.220/30 is directly connected, Serial7/0
R       172.84.18.224/30 [120/2] via 172.84.18.222, 00:00:26, Serial7/0
                         [120/2] via 172.84.18.213, 00:00:03, Serial6/0
                         [120/2] via 172.84.18.201, 00:00:01, Serial3/0
\end{minted}
\caption{Tabla de rutas del router R A4.}
\label{fig:2.1.1.1}
\end{figure}

En cuanto a los aspectos más relevantes:
\begin{enumerate}
    \item Las entradas que aparecen con una \texttt{C} son las subredes que están directamente conectadas al router.
    \item Entre corchetes, aparece en primer lugar la distancia administrativa correspondiente a RIP (120) y en segundo lugar la distacia a la que se encuentra dicha subred, que en el caso de RIP es el número de saltos necesarios.
    \item Por último, nos indica por cuál de nuestras interfaces debemos salir y la IP de nuestro próximo salto.
\end{enumerate}

$\hfill\square$

\vspace{12px}

    \item \textit{Utilizando información de las tablas de rutas y capturas del tráfico RIP
en la red,
explique el funcionamiento de} split horizon \textit{sobre algún enlace de la
red.}

\vspace{7px}

\textsc{Contestación}: Fijándonos en las capturas del tráfico RIP en Figura \ref{fig:2.1.1.2}, observamos que la primera actualización que envía R4 (por la interfaz Serial2/0) no contiene información relevante a la subred 172.84.16.0/23, precisamente porque el router que conecta con dicha interfaz es el que ha proporcionado la información. Las líneas 4 y 5 de la captura manifiestan este hecho.

\begin{figure}[h]
\begin{minted}[linenos, bgcolor=backcolour, frame=single, breaklines=true]{text}
Ejecutado en R 4
Router#debug ip rip
RIP protocol debugging is on
RIP: received v2 update from 172.84.18.218 on Serial2/0
      172.84.16.0/23 via 0.0.0.0 in 1 hops
RIP: sending  v2 update to 224.0.0.9 via Serial2/0 (172.84.18.217)
RIP: build update entries
      172.84.18.0/25 via 0.0.0.0, metric 2, tag 0
      172.84.18.128/26 via 0.0.0.0, metric 2, tag 0
      172.84.18.192/30 via 0.0.0.0, metric 2, tag 0
      172.84.18.196/30 via 0.0.0.0, metric 2, tag 0
      172.84.18.200/30 via 0.0.0.0, metric 1, tag 0
      172.84.18.204/30 via 0.0.0.0, metric 2, tag 0
      172.84.18.208/30 via 0.0.0.0, metric 2, tag 0
      172.84.18.212/30 via 0.0.0.0, metric 1, tag 0
      172.84.18.220/30 via 0.0.0.0, metric 1, tag 0
      172.84.18.224/30 via 0.0.0.0, metric 3, tag 0
\end{minted}
\caption{Tabla de rutas del router R A4.}
\label{fig:2.1.1.2}
\end{figure}

$\hfill\square$

\vspace{12px}

\item \textit{Empleando el comando} tracert \textit{muestre la ruta que sigue el tráfico
desde el Host A5 hasta la
interfaz de R A2 que conecta con la Organización B.}

\vspace{7px}

\textsc{Contestación}: 

\begin{figure}[h]
\begin{minted}[linenos, bgcolor=backcolour, frame=single, breaklines=true]{text}
C:\>tracert 172.84.18.225

Tracing route to 172.84.18.225 over a maximum of 30 hops: 

  1   0 ms      0 ms      0 ms      172.84.17.254
  2   0 ms      0 ms      0 ms      172.84.18.217
  3   0 ms      28 ms     32 ms     172.84.18.201
  4   1 ms      27 ms     31 ms     172.84.18.225

Trace complete.
\end{minted}
\caption{Tabla de rutas del router R A4.}
\label{fig:2.1.1.3}
\end{figure}

$\hfill\square$

\vspace{12px}

\item \textit{Con la simulación en marcha, desactive en R A4 la interfaz de salida
hacia R A2. Utilizando información de las tablas de rutas y capturas
del tráfico RIP en la red, explique en detalle cómo RIP converge a una
nueva solución para alcanzar R A2. Céntrese únicamente en los
routers R A4 y R A2.}

\textit{Indique, en caso de que aplique, el funcionamiento sobre este
escenario y el uso de las técnicas} triggered updates \textit{y} poison reverse.


\vspace{7px}

\textsc{Contestación}: 

Esta es la tabla de rutas de R A4 (Figura \ref{fig:2.1.1.4.1}) y R A2 (Figura \ref{fig:2.1.1.4.2}) antes de desactivar:

\begin{figure}[h]
\begin{minted}[linenos, bgcolor=backcolour, frame=single, breaklines=true]{text}
En R A4

Router>en
Router#sh ip ro
Router#sh ip route 
Codes: C - connected, S - static, I - IGRP, R - RIP, M - mobile, B - BGP
       D - EIGRP, EX - EIGRP external, O - OSPF, IA - OSPF inter area
       N1 - OSPF NSSA external type 1, N2 - OSPF NSSA external type 2
       E1 - OSPF external type 1, E2 - OSPF external type 2, E - EGP
       i - IS-IS, L1 - IS-IS level-1, L2 - IS-IS level-2, ia - IS-IS inter area
       * - candidate default, U - per-user static route, o - ODR
       P - periodic downloaded static route

Gateway of last resort is not set

     172.84.0.0/16 is variably subnetted, 12 subnets, 4 masks
R       172.84.16.0/23 [120/1] via 172.84.18.218, 00:00:08, Serial2/0
R       172.84.18.0/25 [120/1] via 172.84.18.201, 00:00:08, Serial3/0
R       172.84.18.128/26 [120/1] via 172.84.18.222, 00:00:08, Serial7/0
R       172.84.18.192/30 [120/1] via 172.84.18.201, 00:00:08, Serial3/0
R       172.84.18.196/30 [120/1] via 172.84.18.213, 00:00:08, Serial6/0
                         [120/1] via 172.84.18.201, 00:00:08, Serial3/0
C       172.84.18.200/30 is directly connected, Serial3/0
R       172.84.18.204/30 [120/1] via 172.84.18.213, 00:00:08, Serial6/0
R       172.84.18.208/30 [120/1] via 172.84.18.222, 00:00:08, Serial7/0
C       172.84.18.212/30 is directly connected, Serial6/0
C       172.84.18.216/30 is directly connected, Serial2/0
C       172.84.18.220/30 is directly connected, Serial7/0
R       172.84.18.224/30 [120/2] via 172.84.18.222, 00:00:08, Serial7/0
                         [120/2] via 172.84.18.201, 00:00:08, Serial3/0
                         [120/2] via 172.84.18.213, 00:00:08, Serial6/0
\end{minted}
\caption{Tabla de rutas del router R A4.}
\label{fig:2.1.1.4.1}
\end{figure}

\begin{figure}[h]
\begin{minted}[linenos, bgcolor=backcolour, frame=single, breaklines=true]{text}
En R A2

Router>en
Router#sh ip route
Codes: C - connected, S - static, I - IGRP, R - RIP, M - mobile, B - BGP
       D - EIGRP, EX - EIGRP external, O - OSPF, IA - OSPF inter area
       N1 - OSPF NSSA external type 1, N2 - OSPF NSSA external type 2
       E1 - OSPF external type 1, E2 - OSPF external type 2, E - EGP
       i - IS-IS, L1 - IS-IS level-1, L2 - IS-IS level-2, ia - IS-IS inter area
       * - candidate default, U - per-user static route, o - ODR
       P - periodic downloaded static route

Gateway of last resort is not set

     172.84.0.0/16 is variably subnetted, 12 subnets, 4 masks
R       172.84.16.0/23 [120/3] via 172.84.18.210, 00:00:12, Serial3/0
                       [120/3] via 172.84.18.206, 00:00:14, Serial6/0
                       [120/3] via 172.84.18.193, 00:00:14, Serial2/0
R       172.84.18.0/25 [120/1] via 172.84.18.193, 00:00:14, Serial2/0
R       172.84.18.128/26 [120/1] via 172.84.18.210, 00:00:12, Serial3/0
C       172.84.18.192/30 is directly connected, Serial2/0
R       172.84.18.196/30 [120/1] via 172.84.18.193, 00:00:14, Serial2/0
                         [120/1] via 172.84.18.206, 00:00:14, Serial6/0
R       172.84.18.200/30 [120/1] via 172.84.18.193, 00:00:14, Serial2/0
C       172.84.18.204/30 is directly connected, Serial6/0
C       172.84.18.208/30 is directly connected, Serial3/0
R       172.84.18.212/30 [120/1] via 172.84.18.206, 00:00:14, Serial6/0
R       172.84.18.216/30 [120/2] via 172.84.18.210, 00:00:12, Serial3/0
                         [120/2] via 172.84.18.206, 00:00:14, Serial6/0
                         [120/2] via 172.84.18.193, 00:00:14, Serial2/0
R       172.84.18.220/30 [120/1] via 172.84.18.210, 00:00:12, Serial3/0
C       172.84.18.224/30 is directly connected, Serial9/0
\end{minted}
\caption{Tabla de rutas del router R A4.}
\label{fig:2.1.1.4.2}
\end{figure}

Ahora en la interfaz \texttt{Serial3/0} de R A4, que es la que va a R A1, por donde ha pasado el mensaje anterior para ir a R A2, ejecuto \texttt{shutdown}, y veo que ya no aparece la subred \texttt{172.84.18.200/30}, que es la los conecta.



\begin{mdframed}
\begin{verbatim}
En R A4

Router#conf t
Enter configuration commands, one per line.  End with CNTL/Z.
Router(config)#in
Router(config)#interface se
Router(config)#interface serial 3/0
Router(config-if)#shu
Router(config-if)#shutdown 

Router(config-if)#exit
Router(config)#exit
Router#sh ip ro
Router#sh ip route 
Codes: C - connected, S - static, I - IGRP, R - RIP, M - mobile, B - BGP
       D - EIGRP, EX - EIGRP external, O - OSPF, IA - OSPF inter area
       N1 - OSPF NSSA external type 1, N2 - OSPF NSSA external type 2
       E1 - OSPF external type 1, E2 - OSPF external type 2, E - EGP
       i - IS-IS, L1 - IS-IS level-1, L2 - IS-IS level-2, ia - IS-IS inter area
       * - candidate default, U - per-user static route, o - ODR
       P - periodic downloaded static route

Gateway of last resort is not set

     172.84.0.0/16 is variably subnetted, 9 subnets, 3 masks
R       172.84.16.0/23 [120/1] via 172.84.18.218, 00:00:09, Serial2/0
R       172.84.18.128/26 [120/1] via 172.84.18.222, 00:00:03, Serial7/0
R       172.84.18.196/30 [120/1] via 172.84.18.213, 00:00:09, Serial6/0
R       172.84.18.204/30 [120/1] via 172.84.18.213, 00:00:09, Serial6/0
R       172.84.18.208/30 [120/1] via 172.84.18.222, 00:00:03, Serial7/0
C       172.84.18.212/30 is directly connected, Serial6/0
C       172.84.18.216/30 is directly connected, Serial2/0
C       172.84.18.220/30 is directly connected, Serial7/0
R       172.84.18.224/30 [120/2] via 172.84.18.222, 00:00:03, Serial7/0
                         [120/2] via 172.84.18.213, 00:00:09, Serial6/0
\end{verbatim}
\end{mdframed}

Como a R A2 todavía no le ha llegado la información, sigue mostranto que para ir a \texttt{172.84.18.216/30} (la que conecta R A4 y R A5) puede seguir pasando por \texttt{172.84.18.193}, que es R A1, que seguirá siendo uno de los caminos más cortos.

\begin{mdframed}
\begin{verbatim}
En R A2

Router>en
Router#sh ip ro
Router#sh ip route 
Codes: C - connected, S - static, I - IGRP, R - RIP, M - mobile, B - BGP
       D - EIGRP, EX - EIGRP external, O - OSPF, IA - OSPF inter area
       N1 - OSPF NSSA external type 1, N2 - OSPF NSSA external type 2
       E1 - OSPF external type 1, E2 - OSPF external type 2, E - EGP
       i - IS-IS, L1 - IS-IS level-1, L2 - IS-IS level-2, ia - IS-IS inter area
       * - candidate default, U - per-user static route, o - ODR
       P - periodic downloaded static route

Gateway of last resort is not set

     172.84.0.0/16 is variably subnetted, 12 subnets, 4 masks
R       172.84.16.0/23 [120/3] via 172.84.18.210, 00:00:03, Serial3/0
                       [120/3] via 172.84.18.193, 00:00:17, Serial2/0
                       [120/3] via 172.84.18.206, 00:00:09, Serial6/0
R       172.84.18.0/25 [120/1] via 172.84.18.193, 00:00:17, Serial2/0
R       172.84.18.128/26 [120/1] via 172.84.18.210, 00:00:03, Serial3/0
C       172.84.18.192/30 is directly connected, Serial2/0
R       172.84.18.196/30 [120/1] via 172.84.18.193, 00:00:17, Serial2/0
                         [120/1] via 172.84.18.206, 00:00:09, Serial6/0
R       172.84.18.200/30 [120/1] via 172.84.18.193, 00:00:17, Serial2/0
C       172.84.18.204/30 is directly connected, Serial6/0
C       172.84.18.208/30 is directly connected, Serial3/0
R       172.84.18.212/30 [120/1] via 172.84.18.206, 00:00:09, Serial6/0
R       172.84.18.216/30 [120/2] via 172.84.18.210, 00:00:03, Serial3/0
                         [120/2] via 172.84.18.193, 00:00:17, Serial2/0
                         [120/2] via 172.84.18.206, 00:00:09, Serial6/0
R       172.84.18.220/30 [120/1] via 172.84.18.210, 00:00:03, Serial3/0
C       172.84.18.224/30 is directly connected, Serial9/0
\end{verbatim}
\end{mdframed}

Al ejecutar el \texttt{shutdown}, se intercambiarán estos paquetes (en la primera imagen son los que se envían, y en la segunda los que se reciben)

\begin{figure}[H]
    \centering
    \includegraphics[width=0.5\linewidth]{a1.png}
    \label{fig:placeholder}
\end{figure}

\begin{figure}[H]
    \centering
    \includegraphics[width=0.5\linewidth]{a2.png}
    \label{fig:placeholder}
\end{figure}

Aquí podemos observar el \texttt{triggered update}, ya que al instante de desactivar la interfaz, se envían paquetes de actualización.

Asismismo, si nos metemos dentro de los mensajes, por ejemplo recibido por R A3, podemos apreciar la técnica de \texttt{poison reverse}, ya que ha recibido de R A4 que para ir a la subred \texttt{172.84.18.0/25} su coste es 16 (infinito en RIP), dado que la ruta más corta para llegar a dicha subred para R A4 debe pasar por R A3

\begin{figure}[H]
    \centering
    \includegraphics[width=0.5\linewidth]{b1.png}
    \label{fig:placeholder}
\end{figure}

\begin{figure}[H]
    \centering
    \includegraphics[width=0.5\linewidth]{b2.png}
    \label{fig:placeholder}
\end{figure}

R A2 ha recibido un paquete de R A1, y su tabla de rutas va a cambiar. Concretamente se aprecia que para ir a \texttt{172.84.18.216/30}, que es la interfaz de R A4 que conecta con R A5, se le ha eliminado la posibilidad de ir por \texttt{172.84.18.193}, que es R A1

\begin{mdframed}
\begin{verbatim}
En R A2

Router#sh ip ro
Router#sh ip route 
Codes: C - connected, S - static, I - IGRP, R - RIP, M - mobile, B - BGP
       D - EIGRP, EX - EIGRP external, O - OSPF, IA - OSPF inter area
       N1 - OSPF NSSA external type 1, N2 - OSPF NSSA external type 2
       E1 - OSPF external type 1, E2 - OSPF external type 2, E - EGP
       i - IS-IS, L1 - IS-IS level-1, L2 - IS-IS level-2, ia - IS-IS inter area
       * - candidate default, U - per-user static route, o - ODR
       P - periodic downloaded static route

Gateway of last resort is not set

     172.84.0.0/16 is variably subnetted, 12 subnets, 4 masks
R       172.84.16.0/23 [120/3] via 172.84.18.210, 00:00:12, Serial3/0
                       [120/3] via 172.84.18.206, 00:00:14, Serial6/0
R       172.84.18.0/25 [120/1] via 172.84.18.193, 00:00:14, Serial2/0
R       172.84.18.128/26 [120/1] via 172.84.18.210, 00:00:12, Serial3/0
C       172.84.18.192/30 is directly connected, Serial2/0
R       172.84.18.196/30 [120/1] via 172.84.18.193, 00:00:14, Serial2/0
                         [120/1] via 172.84.18.206, 00:00:14, Serial6/0
R       172.84.18.200/30 is possibly down, routing via 172.84.18.193, Serial2/0
C       172.84.18.204/30 is directly connected, Serial6/0
C       172.84.18.208/30 is directly connected, Serial3/0
R       172.84.18.212/30 [120/1] via 172.84.18.206, 00:00:14, Serial6/0
R       172.84.18.216/30 [120/2] via 172.84.18.210, 00:00:12, Serial3/0
                         [120/2] via 172.84.18.206, 00:00:14, Serial6/0
R       172.84.18.220/30 [120/1] via 172.84.18.210, 00:00:12, Serial3/0
C       172.84.18.224/30 is directly connected, Serial9/0
\end{verbatim}
\end{mdframed}

Si seguimos con la simulación

\begin{figure}[H]
    \centering
    \includegraphics[width=0.5\linewidth]{a3.png}
    \label{fig:placeholder}
\end{figure}

Y si vemos las tablas de rutas, vemos que no han cambiado

Si seguimos con la simulación

\begin{figure}[H]
    \centering
    \includegraphics[width=0.5\linewidth]{a4.png}
    \label{fig:placeholder}
\end{figure}

Ya se han terminado de intercambiar mensajes, es decir, ya ha convergido de nuevo la red y las tablas de rutas se quedan como:

\begin{mdframed}
\begin{verbatim}
En R A2

Router#sh ip route 
Codes: C - connected, S - static, I - IGRP, R - RIP, M - mobile, B - BGP
       D - EIGRP, EX - EIGRP external, O - OSPF, IA - OSPF inter area
       N1 - OSPF NSSA external type 1, N2 - OSPF NSSA external type 2
       E1 - OSPF external type 1, E2 - OSPF external type 2, E - EGP
       i - IS-IS, L1 - IS-IS level-1, L2 - IS-IS level-2, ia - IS-IS inter area
       * - candidate default, U - per-user static route, o - ODR
       P - periodic downloaded static route

Gateway of last resort is not set

     172.84.0.0/16 is variably subnetted, 12 subnets, 4 masks
R       172.84.16.0/23 [120/3] via 172.84.18.210, 00:00:13, Serial3/0
                       [120/3] via 172.84.18.206, 00:00:15, Serial6/0
R       172.84.18.0/25 [120/1] via 172.84.18.193, 00:00:15, Serial2/0
R       172.84.18.128/26 [120/1] via 172.84.18.210, 00:00:13, Serial3/0
C       172.84.18.192/30 is directly connected, Serial2/0
R       172.84.18.196/30 [120/1] via 172.84.18.193, 00:00:15, Serial2/0
                         [120/1] via 172.84.18.206, 00:00:15, Serial6/0
R       172.84.18.200/30 is possibly down, routing via 172.84.18.193, Serial2/0
C       172.84.18.204/30 is directly connected, Serial6/0
C       172.84.18.208/30 is directly connected, Serial3/0
R       172.84.18.212/30 [120/1] via 172.84.18.206, 00:00:15, Serial6/0
R       172.84.18.216/30 [120/2] via 172.84.18.210, 00:00:13, Serial3/0
                         [120/2] via 172.84.18.206, 00:00:15, Serial6/0
R       172.84.18.220/30 [120/1] via 172.84.18.210, 00:00:13, Serial3/0
C       172.84.18.224/30 is directly connected, Serial9/0
\end{verbatim}
\end{mdframed}

\begin{mdframed}
\begin{verbatim}
En R A4

Router#sh ip route
Codes: C - connected, S - static, I - IGRP, R - RIP, M - mobile, B - BGP
       D - EIGRP, EX - EIGRP external, O - OSPF, IA - OSPF inter area
       N1 - OSPF NSSA external type 1, N2 - OSPF NSSA external type 2
       E1 - OSPF external type 1, E2 - OSPF external type 2, E - EGP
       i - IS-IS, L1 - IS-IS level-1, L2 - IS-IS level-2, ia - IS-IS inter area
       * - candidate default, U - per-user static route, o - ODR
       P - periodic downloaded static route

Gateway of last resort is not set

     172.84.0.0/16 is variably subnetted, 9 subnets, 3 masks
R       172.84.16.0/23 [120/1] via 172.84.18.218, 00:00:11, Serial2/0
R       172.84.18.128/26 [120/1] via 172.84.18.222, 00:00:13, Serial7/0
R       172.84.18.196/30 [120/1] via 172.84.18.213, 00:00:15, Serial6/0
R       172.84.18.204/30 [120/1] via 172.84.18.213, 00:00:15, Serial6/0
R       172.84.18.208/30 [120/1] via 172.84.18.222, 00:00:13, Serial7/0
C       172.84.18.212/30 is directly connected, Serial6/0
C       172.84.18.216/30 is directly connected, Serial2/0
C       172.84.18.220/30 is directly connected, Serial7/0
R       172.84.18.224/30 [120/2] via 172.84.18.222, 00:00:13, Serial7/0
                         [120/2] via 172.84.18.213, 00:00:15, Serial6/0

\end{verbatim}
\end{mdframed}


$\hfill\square$

\end{enumerate}

\newpage

\subsubsection{Organización B}

\begin{enumerate}

    \item \textit{Realice un traceroute entre R B4 y el H B2. Justifica el camino que siguen los paquetes.}

\vspace{7px}

\textsc{Contestación}: Tal y como refleja el comando \textit{traceroute} ejecutado y expuesto a continuación, para llegar a H B2 cuya IP es 171.84.16.1 desde R B4, se sigue el siguiente camino:
\begin{itemize}
    \item 171.84.16.189 (R B4) $\xrightarrow[\text{Fa1/0}]{}$ 171.84.16.190 (R B3)

    \item 171.84.17.138 (R B3) $\xrightarrow[\text{Se3/0}]{}$ 171.84.17.137 (R B2)

    \item 171.84.16.126 (R B2) $\xrightarrow[\text{Fa0/0}]{}$ 171.84.16.1 (H B2)
\end{itemize}

\begin{mdframed}
\begin{verbatim}
Router>traceroute 171.84.16.1
Type escape sequence to abort.
Tracing the route to 171.84.16.1

  1   171.84.16.190   0 msec    0 msec    0 msec    
  2   171.84.17.137   20 msec   0 msec    26 msec   
  3   171.84.16.1     0 msec    1 msec    0 msec    
\end{verbatim}
\end{mdframed}

$\hfill\square$

\vspace{12px}

    \item \textit{Deshabilite ahora la interfaz del router R B3 que conecta con el R B2. Espere a que la red converja de nuevo. A continuación, realiza el traceroute de nuevo entre R B4 y H B2 y justifica el camino que ahora siguen los paquetes. Cuando termine este ejercicio, vuelta a activar la interfaz desactivada.}

\vspace{7px}

\textsc{Contestación}: 
Para desactivarlo, basta ejecutar en la interfaz \texttt{Serial3/0} de R B3 \texttt{shutdown}

Si ejecutamos ahora \textit{traceroute}, nos mostrará el siguiente camino:
\begin{itemize}
    \item 171.84.16.189 (R B4) $\xrightarrow[\text{Fa1/0}]{}$ 171.84.16.190 (R B3)

    \item 171.84.17.134 (R B3) $\xrightarrow[\text{Se2/0}]{}$ 171.84.17.133 (R B1)

    \item 171.84.17.129 (R B1) $\xrightarrow[\text{Se2/0}]{}$ 171.84.17.130 (R B2)
    
    \item 171.84.16.126 (R B2) $\xrightarrow[\text{Fa0/0}]{}$ 171.84.16.1 (H B2)
\end{itemize}

\begin{mdframed}
\begin{verbatim}
Router>traceroute 171.84.16.1
Type escape sequence to abort.
Tracing the route to 171.84.16.1

  1   171.84.16.190   0 msec    0 msec    0 msec    
  2   171.84.17.133   0 msec    19 msec   0 msec    
  3   171.84.17.130   1 msec    1 msec    0 msec    
  4   171.84.16.1     1 msec    0 msec    1 msec    
\end{verbatim}
\end{mdframed}

$\hfill\square$

\vspace{12px}

\item \textit{Realice la configuración necesaria para que R B4 y R B3 se
conviertan en Designated Router (DR) y Backup Designated Router (BDR), respectivamente, de la LAN B3. Muestre que efectivamente se ha realizado el cambio.}

\vspace{7px}

\textsc{Contestación}: Para hacerlo, nos tendremos que meter en modo de configuración de la interfaz Fa0/0 de R B3 y ejecutar \textit{ip ospf priority 3}, y en el modo de configuración de la interfaz Fa1/0 de R B4 \textit{ip ospf priority 2}. Para visualizar el cambio podemos ejecutar (por ejemplo, en R B4):

\begin{mdframed}
\begin{verbatim}
Router#show ip ospf interface 

...
FastEthernet1/0 is up, line protocol is up
  Internet address is 171.84.16.189/26, Area 2
  Process ID 100, Router ID 171.84.16.254, Network Type BROADCAST, Cost: 1
  Transmit Delay is 1 sec, State DR, Priority 3
  Designated Router (ID) 171.84.16.254, Interface address 171.84.16.189
  Backup Designated Router (ID) 171.84.17.138, Interface address 171.84.16.190
  Timer intervals configured, Hello 10, Dead 40, Wait 40, Retransmit 5
    Hello due in 00:00:05
  Index 2/2, flood queue length 0
  Next 0x0(0)/0x0(0)
  Last flood scan length is 1, maximum is 1
  Last flood scan time is 0 msec, maximum is 0 msec
  Neighbor Count is 1, Adjacent neighbor count is 1
    Adjacent with neighbor 171.84.17.138  (Backup Designated Router)
  Suppress hello for 0 neighbor(s)
\end{verbatim}
\end{mdframed}

Viendo el identificador, podemos ver que el de DR coincide con el de R B4 y el de BDR con el de R B3 (sus identificadores son la mayor dirección IP de sus interfaces). Podemos ver sus identificadores con: 

\begin{mdframed}
\begin{verbatim}
En R B4

Router#show ip protocols 

Routing Protocol is "ospf 100"
  Outgoing update filter list for all interfaces is not set 
  Incoming update filter list for all interfaces is not set 
  Router ID 171.84.16.254
  Number of areas in this router is 1. 1 normal 0 stub 0 nssa
  Maximum path: 4
  Routing for Networks:
    171.84.16.128 0.0.0.63 area 2
    171.84.16.192 0.0.0.63 area 2
  Passive Interface(s): 
    FastEthernet0/0
  Routing Information Sources:  
    Gateway         Distance      Last Update 
    171.84.16.254        110      00:11:51
    171.84.17.138        110      00:11:51
  Distance: (default is 110)
\end{verbatim}
\end{mdframed}

\begin{mdframed}
\begin{verbatim}
En R B3

Router#sh ip protocols 

Routing Protocol is "ospf 100"
  Outgoing update filter list for all interfaces is not set 
  Incoming update filter list for all interfaces is not set 
  Router ID 171.84.17.138
  Number of areas in this router is 3. 3 normal 0 stub 0 nssa
  Maximum path: 4
  Routing for Networks:
    171.84.16.128 0.0.0.63 area 2
    171.84.17.80 0.0.0.3 area 3
    171.84.17.84 0.0.0.3 area 3
    171.84.17.136 0.0.0.3 area 0
    171.84.17.132 0.0.0.3 area 0
  Routing Information Sources:  
    Gateway         Distance      Last Update 
    171.84.16.254        110      00:13:51
    171.84.17.89         110      00:14:26
    171.84.17.90         110      00:14:21
    171.84.17.137        110      00:14:25
    171.84.17.138        110      00:14:26
    172.84.18.226        110      00:14:25
  Distance: (default is 110)
\end{verbatim}
\end{mdframed}

$\hfill\square$

\vspace{12px}

\item \textit{Muestre las tablas de rutas de R B6 y comente los aspectos más relevantes.}

\vspace{7px}

\textsc{Contestación}: 

\begin{mdframed}
\begin{verbatim}
Router#show ip route 
Codes: C - connected, S - static, I - IGRP, R - RIP, M - mobile, B - BGP
       D - EIGRP, EX - EIGRP external, O - OSPF, IA - OSPF inter area
       N1 - OSPF NSSA external type 1, N2 - OSPF NSSA external type 2
       E1 - OSPF external type 1, E2 - OSPF external type 2, E - EGP
       i - IS-IS, L1 - IS-IS level-1, L2 - IS-IS level-2, ia - IS-IS inter area
       * - candidate default, U - per-user static route, o - ODR
       P - periodic downloaded static route

Gateway of last resort is not set

     171.84.0.0/16 is variably subnetted, 11 subnets, 4 masks
O IA    171.84.16.0/25 [110/129] via 171.84.17.85, 00:01:16, Serial2/0
O IA    171.84.16.128/26 [110/65] via 171.84.17.85, 00:01:16, Serial2/0
O IA    171.84.16.192/26 [110/66] via 171.84.17.85, 00:01:16, Serial2/0
C       171.84.17.0/26 is directly connected, FastEthernet0/0
O       171.84.17.64/28 [110/65] via 171.84.17.89, 00:29:16, Serial3/0
O       171.84.17.80/30 [110/128] via 171.84.17.85, 00:01:16, Serial2/0
                        [110/128] via 171.84.17.89, 00:01:16, Serial3/0
C       171.84.17.84/30 is directly connected, Serial2/0
C       171.84.17.88/30 is directly connected, Serial3/0
O IA    171.84.17.128/30 [110/192] via 171.84.17.85, 00:01:16, Serial2/0
O IA    171.84.17.132/30 [110/128] via 171.84.17.85, 00:01:16, Serial2/0
O IA    171.84.17.136/30 [110/128] via 171.84.17.85, 00:01:16, Serial2/0
\end{verbatim}
\end{mdframed}

Podemos destacar varios aspectos: 

\begin{itemize}
    \item[-] Lo que va entre corchetes son, en primer lugar, la distancia administrativa del protocolo de enrutamiento usado (en este caso 110, OSPF), y el segundo lugar es el coste para llegar a la red destino \[\text{Coste} = \frac{10^8}{\text{ancho de banda (bps)}}\]
    \item[-] Los destinos que aparecen con una \textit{C} significan que están directamente conectados a mi router, y que se llegan a ellos por la interfaz que aparezca.
    \item[-] Las que se muestran solamente con una \textit{O} indican las rutas aprendidas por OSPF y que pertenecen a alguna de las áreas en la que está el router, y se llegan a ellas por la IP que aparece, que está conectado a la interfaz de la router que me indica.
    \item[-] Las que tienen una \textit{O IA} son como las anteriores, pero que se encuentran en una área distinta a las del router.
    \item[-] Por último, para la subred 171.84.17.80/30 aparecen dos líneas, esto quiere decir que se puede llegar a ella por dos distintos caminos, con el mismo coste para ambos.
\end{itemize}


$\hfill\square$

\vspace{12px}

\item \textit{Realice la configuración necesaria para que el área 3 sea una stub
area. Analizando las tablas de rutas que considere relevantes, ¿qué
diferencias observa con respecto a la configuración anterior? ¿Por
qué?}

\vspace{7px}

\textsc{Contestación}: 
Para configurarlo solo es necesario ejecutar (en el modo de configuración de OSPF) en todos los routers del área (en este caso R B3, R B5, R B6) \texttt{area 3 stub}

La diferencia está en que se ha añadido una ruta por defecto para todos aquellos paquetes que vaya a una red externa (fuera del AS) para los routers R B5 y R B6 (los que no son ABR). Para estos paquetes, estos deben pasar por el ABR, ya que esta área carece de un ASBR, por lo que todos los que vayan a una red externa deben pasar por R B3 (que es el único ABR del área 3).
\begin{mdframed}
\begin{verbatim}
En R B5

Router#sh ip route 
Codes: C - connected, S - static, I - IGRP, R - RIP, M - mobile, B - BGP
       D - EIGRP, EX - EIGRP external, O - OSPF, IA - OSPF inter area
       N1 - OSPF NSSA external type 1, N2 - OSPF NSSA external type 2
       E1 - OSPF external type 1, E2 - OSPF external type 2, E - EGP
       i - IS-IS, L1 - IS-IS level-1, L2 - IS-IS level-2, ia - IS-IS inter area
       * - candidate default, U - per-user static route, o - ODR
       P - periodic downloaded static route

Gateway of last resort is 171.84.17.81 to network 0.0.0.0

     171.84.0.0/16 is variably subnetted, 11 subnets, 4 masks
O IA    171.84.16.0/25 [110/129] via 171.84.17.81, 00:01:44, Serial2/0
O IA    171.84.16.128/26 [110/65] via 171.84.17.81, 00:01:44, Serial2/0
O IA    171.84.16.192/26 [110/66] via 171.84.17.81, 00:01:44, Serial2/0
O       171.84.17.0/26 [110/65] via 171.84.17.90, 00:02:09, Serial3/0
C       171.84.17.64/28 is directly connected, FastEthernet0/0
C       171.84.17.80/30 is directly connected, Serial2/0
O       171.84.17.84/30 [110/128] via 171.84.17.90, 00:01:44, Serial3/0
                        [110/128] via 171.84.17.81, 00:01:44, Serial2/0
C       171.84.17.88/30 is directly connected, Serial3/0
O IA    171.84.17.128/30 [110/192] via 171.84.17.81, 00:01:44, Serial2/0
O IA    171.84.17.132/30 [110/128] via 171.84.17.81, 00:01:44, Serial2/0
O IA    171.84.17.136/30 [110/128] via 171.84.17.81, 00:01:44, Serial2/0
O*IA 0.0.0.0/0 [110/65] via 171.84.17.81, 00:01:44, Serial2/0
\end{verbatim}
\end{mdframed}

\begin{mdframed}
\begin{verbatim}
En R B6

Router#sh ip route 
Codes: C - connected, S - static, I - IGRP, R - RIP, M - mobile, B - BGP
       D - EIGRP, EX - EIGRP external, O - OSPF, IA - OSPF inter area
       N1 - OSPF NSSA external type 1, N2 - OSPF NSSA external type 2
       E1 - OSPF external type 1, E2 - OSPF external type 2, E - EGP
       i - IS-IS, L1 - IS-IS level-1, L2 - IS-IS level-2, ia - IS-IS inter area
       * - candidate default, U - per-user static route, o - ODR
       P - periodic downloaded static route

Gateway of last resort is 171.84.17.85 to network 0.0.0.0

     171.84.0.0/16 is variably subnetted, 11 subnets, 4 masks
O IA    171.84.16.0/25 [110/129] via 171.84.17.85, 00:01:49, Serial2/0
O IA    171.84.16.128/26 [110/65] via 171.84.17.85, 00:01:49, Serial2/0
O IA    171.84.16.192/26 [110/66] via 171.84.17.85, 00:01:49, Serial2/0
C       171.84.17.0/26 is directly connected, FastEthernet0/0
O       171.84.17.64/28 [110/65] via 171.84.17.89, 00:02:15, Serial3/0
O       171.84.17.80/30 [110/128] via 171.84.17.85, 00:01:49, Serial2/0
                        [110/128] via 171.84.17.89, 00:01:49, Serial3/0
C       171.84.17.84/30 is directly connected, Serial2/0
C       171.84.17.88/30 is directly connected, Serial3/0
O IA    171.84.17.128/30 [110/192] via 171.84.17.85, 00:01:49, Serial2/0
O IA    171.84.17.132/30 [110/128] via 171.84.17.85, 00:01:49, Serial2/0
O IA    171.84.17.136/30 [110/128] via 171.84.17.85, 00:01:49, Serial2/0
O*IA 0.0.0.0/0 [110/65] via 171.84.17.85, 00:01:49, Serial2/0
\end{verbatim}
\end{mdframed}

$\hfill\square$

\vspace{12px}

\item \textit{Muestre las tablas de rutas de R B4 y comente los aspectos más
relevantes. ¿Cuál es el camino óptimo para alcanzar R B5?}

\vspace{7px}

\textsc{Contestación}:

\begin{mdframed}
\begin{verbatim}
Router#sh ip route 
Codes: C - connected, S - static, I - IGRP, R - RIP, M - mobile, B - BGP
       D - EIGRP, EX - EIGRP external, O - OSPF, IA - OSPF inter area
       N1 - OSPF NSSA external type 1, N2 - OSPF NSSA external type 2
       E1 - OSPF external type 1, E2 - OSPF external type 2, E - EGP
       i - IS-IS, L1 - IS-IS level-1, L2 - IS-IS level-2, ia - IS-IS inter area
       * - candidate default, U - per-user static route, o - ODR
       P - periodic downloaded static route

Gateway of last resort is not set

     171.84.0.0/16 is variably subnetted, 11 subnets, 4 masks
O IA    171.84.16.0/25 [110/66] via 171.84.16.190, 00:00:00, FastEthernet1/0
C       171.84.16.128/26 is directly connected, FastEthernet1/0
C       171.84.16.192/26 is directly connected, FastEthernet0/0
O IA    171.84.17.0/26 [110/66] via 171.84.16.190, 00:00:00, FastEthernet1/0
O IA    171.84.17.64/28 [110/66] via 171.84.16.190, 00:00:00, FastEthernet1/0
O IA    171.84.17.80/30 [110/65] via 171.84.16.190, 00:00:00, FastEthernet1/0
O IA    171.84.17.84/30 [110/65] via 171.84.16.190, 00:00:00, FastEthernet1/0
O IA    171.84.17.88/30 [110/129] via 171.84.16.190, 00:00:00, FastEthernet1/0
O IA    171.84.17.128/30 [110/129] via 171.84.16.190, 00:00:00, FastEthernet1/0
O IA    171.84.17.132/30 [110/65] via 171.84.16.190, 00:00:00, FastEthernet1/0
O IA    171.84.17.136/30 [110/65] via 171.84.16.190, 00:00:00, FastEthernet1/0
\end{verbatim}
\end{mdframed}

Podemos destacar varios aspectos:

\begin{itemize}
    \item[-] Para cada entrada de la tabla de rutas, los valores representados entre corchetes son, por una parte, la distancia administrativa del protocolo de enrutamiento usado y, por otro lado, el coste para llegar a la red de destino \[\text{Coste} = \frac{10^8}{\text{ancho de banda (bps)}}\]

    \item[-] Los destinos cuya entrada está precedida por una $C$ son aquellos conectados directamente al router en estudio.

    \item[-] Los que están precedidos por $O$ $IA$ indican rutas aprendidas por OSPF y que se encuentran en un área distinta a las del router.
\end{itemize}

En cuanto al interrogante adicional, el camino óptimo se cumple pasando por el R B3.

$\hfill\square$

\vspace{12px}

\item \textit{Realice la configuración necesaria para que el área 2 sea una totally
stub area. Analizando las tablas de rutas que considere relevantes,
demuestre que se trata de una totally stub área. ¿Qué diferencias
observa con respecto a la configuración anterior? ¿Por qué?}


\vspace{7px}

\textsc{Contestación}: 
Para configurarlo, debemos ejecutar (en el modo de configuración de OSPF) en todos los router internos (en este caso solo R B4) \texttt{area 2 stub} y en los ABR (en este caso R B3) \texttt{area 2 stub no-summary}

La diferencia está en que se ha dejado solo las rutas a las que router está directamente conetados, y el resto se han agrupado en una por defecto. Esto se debe a que en una totally stub area, para llegar a cualquier subred que no se encuentra en el área, la única alternativa es pasar por un único ABR (en este caso el ABR es R B3).

\begin{mdframed}
\begin{verbatim}
En R B4

Router#sh ip route 
Codes: C - connected, S - static, I - IGRP, R - RIP, M - mobile, B - BGP
       D - EIGRP, EX - EIGRP external, O - OSPF, IA - OSPF inter area
       N1 - OSPF NSSA external type 1, N2 - OSPF NSSA external type 2
       E1 - OSPF external type 1, E2 - OSPF external type 2, E - EGP
       i - IS-IS, L1 - IS-IS level-1, L2 - IS-IS level-2, ia - IS-IS inter area
       * - candidate default, U - per-user static route, o - ODR
       P - periodic downloaded static route

Gateway of last resort is 171.84.16.190 to network 0.0.0.0

     171.84.0.0/26 is subnetted, 2 subnets
C       171.84.16.128 is directly connected, FastEthernet1/0
C       171.84.16.192 is directly connected, FastEthernet0/0
O*IA 0.0.0.0/0 [110/2] via 171.84.16.190, 00:00:18, FastEthernet1/0
\end{verbatim}
\end{mdframed}

$\hfill\square$

\vspace{12px}

\item \textit{Realice la configuración necesaria (sin desactivar las interfaces) para
que el camino óptimo entre R B4 y H B5 pase a través de R B6.
Muestre que efectivamente se ha realizado el cambio.}

\vspace{7px}

\textsc{Contestación}: 
\begin{mdframed}
\begin{verbatim}
En R B5

Router(config)#interface serial 3/0
Router(config-if)#ip ospf cost 31
\end{verbatim}
\end{mdframed}

\begin{mdframed}
\begin{verbatim}
En R B6

Router(config)#interface serial 2/0
Router(config-if)#ip ospf cost 31
\end{verbatim}
\end{mdframed}

\begin{mdframed}
\begin{verbatim}
En R B5

Router#sh ip route 
Codes: C - connected, S - static, I - IGRP, R - RIP, M - mobile, B - BGP
       D - EIGRP, EX - EIGRP external, O - OSPF, IA - OSPF inter area
       N1 - OSPF NSSA external type 1, N2 - OSPF NSSA external type 2
       E1 - OSPF external type 1, E2 - OSPF external type 2, E - EGP
       i - IS-IS, L1 - IS-IS level-1, L2 - IS-IS level-2, ia - IS-IS inter area
       * - candidate default, U - per-user static route, o - ODR
       P - periodic downloaded static route

Gateway of last resort is 171.84.17.90 to network 0.0.0.0

     171.84.0.0/16 is variably subnetted, 11 subnets, 4 masks
O IA    171.84.16.0/25 [110/127] via 171.84.17.90, 00:08:12, Serial3/0
O IA    171.84.16.128/26 [110/63] via 171.84.17.90, 00:08:12, Serial3/0
O IA    171.84.16.192/26 [110/64] via 171.84.17.90, 00:08:12, Serial3/0
O       171.84.17.0/26 [110/32] via 171.84.17.90, 00:08:42, Serial3/0
C       171.84.17.64/28 is directly connected, FastEthernet0/0
C       171.84.17.80/30 is directly connected, Serial2/0
O       171.84.17.84/30 [110/62] via 171.84.17.90, 00:08:12, Serial3/0
C       171.84.17.88/30 is directly connected, Serial3/0
O IA    171.84.17.128/30 [110/190] via 171.84.17.90, 00:08:12, Serial3/0
O IA    171.84.17.132/30 [110/126] via 171.84.17.90, 00:08:12, Serial3/0
O IA    171.84.17.136/30 [110/126] via 171.84.17.90, 00:08:12, Serial3/0
O*IA 0.0.0.0/0 [110/63] via 171.84.17.90, 00:08:12, Serial3/0
\end{verbatim}
\end{mdframed}

Podemos ver que para llegar ahora a la subred 171.84.17.128/30 en la que se encuenrtra R B4, se debe pasar por 171.84.17.90, que es una interfaz de R B6.

$\hfill\square$

\vspace{12px}

\item \textit{Utilizando la herramienta Cisco Packet Tracer capture tráfico OSPF
para mostrar al menos cuatro tipos de LSAs diferentes que se
intercambian los routers del escenario e indique su propósito.}

\vspace{7px}

\textsc{Contestación}: 
\begin{figure}[H]
    \centering
    \includegraphics[width=0.9\linewidth]{LSA1.1.png}
    \label{fig:placeholder}
\end{figure}

\begin{figure}[H]
    \centering
    \includegraphics[width=0.9\linewidth]{LSA1.2.png}
    \label{fig:placeholder}
\end{figure}

\begin{figure}[H]
    \centering
    \includegraphics[width=0.9\linewidth]{LSA2.png}
    \label{fig:placeholder}
\end{figure}

\begin{figure}[H]
    \centering
    \includegraphics[width=0.5\linewidth]{LSA3.png}
    \label{fig:placeholder}
\end{figure}

$\hfill\square$

\vspace{12px}

\item \textit{Configura ahora los routers ABR de cada área para que realicen la
agregación de las subredes que se encuentran en su área.
Comprueba que se ha realizado esta agregación de redes mostrando
la tabla de rutas del R B2 antes y después de aplicar dicha
agregación.}

\vspace{7px}

\textsc{Contestación}: 

\begin{mdframed}
\begin{verbatim}
R B2 antes:

Router#sh ip route 
Codes: C - connected, S - static, I - IGRP, R - RIP, M - mobile, B - BGP
       D - EIGRP, EX - EIGRP external, O - OSPF, IA - OSPF inter area
       N1 - OSPF NSSA external type 1, N2 - OSPF NSSA external type 2
       E1 - OSPF external type 1, E2 - OSPF external type 2, E - EGP
       i - IS-IS, L1 - IS-IS level-1, L2 - IS-IS level-2, ia - IS-IS inter area
       * - candidate default, U - per-user static route, o - ODR
       P - periodic downloaded static route

Gateway of last resort is not set

     171.84.0.0/16 is variably subnetted, 11 subnets, 4 masks
C       171.84.16.0/25 is directly connected, FastEthernet0/0
O IA    171.84.16.128/26 [110/65] via 171.84.17.138, 00:03:42, Serial3/0
O IA    171.84.16.192/26 [110/66] via 171.84.17.138, 00:03:27, Serial3/0
O IA    171.84.17.0/26 [110/129] via 171.84.17.138, 00:03:42, Serial3/0
O IA    171.84.17.64/28 [110/129] via 171.84.17.138, 00:03:42, Serial3/0
O IA    171.84.17.80/30 [110/128] via 171.84.17.138, 00:03:42, Serial3/0
O IA    171.84.17.84/30 [110/128] via 171.84.17.138, 00:03:42, Serial3/0
O IA    171.84.17.88/30 [110/159] via 171.84.17.138, 00:03:42, Serial3/0
C       171.84.17.128/30 is directly connected, Serial2/0
O       171.84.17.132/30 [110/128] via 171.84.17.138, 00:03:42, Serial3/0
                         [110/128] via 171.84.17.129, 00:03:42, Serial2/0
C       171.84.17.136/30 is directly connected, Serial3/0
\end{verbatim}
\end{mdframed}

\begin{mdframed}
\begin{verbatim}
En R B2:

Router#conf t
Enter configuration commands, one per line.  End with CNTL/Z.
Router(config)#router osp
Router(config)#router ospf 100
Router(config-router)#area 1 range 171.84.16.0 255.255.255.128
\end{verbatim}
\end{mdframed}

\begin{mdframed}
\begin{verbatim}
En R B3:

Router>en
Router#conf t
Enter configuration commands, one per line.  End with CNTL/Z.
Router(config)#router os
Router(config)#router ospf 100
Router(config-router)#area 2 range 171.84.16.128 255.255.255.128
Router(config-router)#area 3 range 171.84.17.0 255.255.255.128
\end{verbatim}
\end{mdframed}

\begin{mdframed}
\begin{verbatim}
En R B2:

Router#sh ip route
Codes: C - connected, S - static, I - IGRP, R - RIP, M - mobile, B - BGP
       D - EIGRP, EX - EIGRP external, O - OSPF, IA - OSPF inter area
       N1 - OSPF NSSA external type 1, N2 - OSPF NSSA external type 2
       E1 - OSPF external type 1, E2 - OSPF external type 2, E - EGP
       i - IS-IS, L1 - IS-IS level-1, L2 - IS-IS level-2, ia - IS-IS inter area
       * - candidate default, U - per-user static route, o - ODR
       P - periodic downloaded static route

Gateway of last resort is not set

     171.84.0.0/16 is variably subnetted, 6 subnets, 2 masks
C       171.84.16.0/25 is directly connected, FastEthernet0/0
O IA    171.84.16.128/25 [110/66] via 171.84.17.138, 00:00:45, Serial3/0
O IA    171.84.17.0/25 [110/128] via 171.84.17.138, 00:00:27, Serial3/0
C       171.84.17.128/30 is directly connected, Serial2/0
O       171.84.17.132/30 [110/128] via 171.84.17.138, 00:10:47, Serial3/0
                         [110/128] via 171.84.17.129, 00:10:47, Serial2/0
C       171.84.17.136/30 is directly connected, Serial3/0
\end{verbatim}
\end{mdframed}


$\hfill\square$

\end{enumerate}

\subsection{Interconexión y Redistribución de Rutas}

\begin{enumerate}
    \item \textit{Muestre las tablas de rutas de R A2 y R B1 y coméntelas en detalle.}
    \vspace{7px}

    \textsc{Contestación}: 
    
\begin{mdframed}
\begin{verbatim}
Para configurar todo, en R B1

Router>en
Router#conf t
Enter configuration commands, one per line.  End with CNTL/Z.
Router(config)#router rip
Router(config-router)#ver
Router(config-router)#version 2
Router(config-router)#network 172.84.18.224
Router(config-router)#pas
Router(config-router)#passive-interface se
Router(config-router)#passive-interface serial 2/0
Router(config-router)#passive-interface serial 3/0
Router(config-router)#redistr
Router(config-router)#redistribute ospf 100 met
Router(config-router)#redistribute ospf 100 metric 1
Router(config-router)#exit
Router(config)#router ospf 100
Router(config-router)#re
Router(config-router)#redistribute rip metric 200 subnet
Router(config-router)#pass
Router(config-router)#passive-interface se
Router(config-router)#passive-interface serial 9/0
\end{verbatim}
\end{mdframed}

En R B1 se muestra:
    \begin{enumerate}
        \item Una primera parte que corresponde a OSPF, que son las que entre corchetes el primer número es 110 y aparecen con un O, con la correspondiente métrica de OSPF a continucación, que depende del ancho de banda. Además, las que llevan IA indica que se encuentra en un área distinto al que pertenece R B1, el área 0
        \item Por otra parte, las que llevan R, cuya distancia administrativa es 120, y el coste es el número de saltos.
        \item También, las que aparecen que un C son las redes directamentes conectadas a R B1
    \end{enumerate}

\begin{mdframed}
\begin{verbatim}
En R B1

Router#sh ip route 
Codes: C - connected, S - static, I - IGRP, R - RIP, M - mobile, B - BGP
       D - EIGRP, EX - EIGRP external, O - OSPF, IA - OSPF inter area
       N1 - OSPF NSSA external type 1, N2 - OSPF NSSA external type 2
       E1 - OSPF external type 1, E2 - OSPF external type 2, E - EGP
       i - IS-IS, L1 - IS-IS level-1, L2 - IS-IS level-2, ia - IS-IS inter area
       * - candidate default, U - per-user static route, o - ODR
       P - periodic downloaded static route

Gateway of last resort is not set

     171.84.0.0/16 is variably subnetted, 7 subnets, 3 masks
R       171.84.0.0/16 [120/10] via 172.84.18.225, 00:02:55, Serial9/0
O IA    171.84.16.0/25 [110/65] via 171.84.17.130, 00:24:12, Serial2/0
O IA    171.84.16.128/25 [110/66] via 171.84.17.134, 00:14:05, Serial3/0
O IA    171.84.17.0/25 [110/128] via 171.84.17.134, 00:13:47, Serial3/0
C       171.84.17.128/30 is directly connected, Serial2/0
C       171.84.17.132/30 is directly connected, Serial3/0
O       171.84.17.136/30 [110/128] via 171.84.17.130, 00:24:12, Serial2/0
                         [110/128] via 171.84.17.134, 00:24:12, Serial3/0
     172.84.0.0/16 is variably subnetted, 12 subnets, 4 masks
R       172.84.16.0/23 [120/4] via 172.84.18.225, 00:00:18, Serial9/0
R       172.84.18.0/25 [120/2] via 172.84.18.225, 00:00:18, Serial9/0
R       172.84.18.128/26 [120/2] via 172.84.18.225, 00:00:18, Serial9/0
R       172.84.18.192/30 [120/1] via 172.84.18.225, 00:00:18, Serial9/0
R       172.84.18.196/30 [120/2] via 172.84.18.225, 00:00:18, Serial9/0
R       172.84.18.200/30 [120/2] via 172.84.18.225, 00:00:18, Serial9/0
R       172.84.18.204/30 [120/1] via 172.84.18.225, 00:00:18, Serial9/0
R       172.84.18.208/30 [120/1] via 172.84.18.225, 00:00:18, Serial9/0
R       172.84.18.212/30 [120/2] via 172.84.18.225, 00:00:18, Serial9/0
R       172.84.18.216/30 [120/3] via 172.84.18.225, 00:00:18, Serial9/0
R       172.84.18.220/30 [120/2] via 172.84.18.225, 00:00:18, Serial9/0
C       172.84.18.224/30 is directly connected, Serial9/0
\end{verbatim}
\end{mdframed}

En R A2 se muestra:

\begin{enumerate}
    \item En primer lugar, cabe destacar que todo está en RIP, ya que R B1 ha redistribuido todas las de OSPF a RIP con métrica 1, que se ve que para ir a 171.84.0.0/16 tiene coste 1.
    \item Otro cambio notable, es que todas las rutas de la organización B aparece como \texttt{is possibly down} porque son rutas que has aprendido por la redistribución de rutas, por lo que no sabes exactamente la situación de esas subredes. 
\end{enumerate}

\begin{mdframed}
\begin{verbatim}
En R A2

Router#sh ip route 
Codes: C - connected, S - static, I - IGRP, R - RIP, M - mobile, B - BGP
       D - EIGRP, EX - EIGRP external, O - OSPF, IA - OSPF inter area
       N1 - OSPF NSSA external type 1, N2 - OSPF NSSA external type 2
       E1 - OSPF external type 1, E2 - OSPF external type 2, E - EGP
       i - IS-IS, L1 - IS-IS level-1, L2 - IS-IS level-2, ia - IS-IS inter area
       * - candidate default, U - per-user static route, o - ODR
       P - periodic downloaded static route

Gateway of last resort is not set

     171.84.0.0/16 is variably subnetted, 7 subnets, 3 masks
R       171.84.0.0/16 [120/1] via 172.84.18.226, 00:00:16, Serial9/0
R       171.84.16.0/25 is possibly down, routing via 172.84.18.226, Serial9/0
R       171.84.16.128/25 is possibly down, routing via 172.84.18.226, Serial9/0
R       171.84.17.0/25 is possibly down, routing via 172.84.18.226, Serial9/0
R       171.84.17.128/30 is possibly down, routing via 172.84.18.226, Serial9/0
R       171.84.17.132/30 is possibly down, routing via 172.84.18.226, Serial9/0
R       171.84.17.136/30 is possibly down, routing via 172.84.18.226, Serial9/0
     172.84.0.0/16 is variably subnetted, 12 subnets, 4 masks
R       172.84.16.0/23 [120/3] via 172.84.18.206, 00:00:04, Serial6/0
                       [120/3] via 172.84.18.210, 00:00:22, Serial3/0
                       [120/3] via 172.84.18.193, 00:00:05, Serial2/0
R       172.84.18.0/25 [120/1] via 172.84.18.193, 00:00:05, Serial2/0
R       172.84.18.128/26 [120/1] via 172.84.18.210, 00:00:22, Serial3/0
C       172.84.18.192/30 is directly connected, Serial2/0
R       172.84.18.196/30 [120/1] via 172.84.18.206, 00:00:04, Serial6/0
                         [120/1] via 172.84.18.193, 00:00:05, Serial2/0
R       172.84.18.200/30 [120/1] via 172.84.18.193, 00:00:05, Serial2/0
C       172.84.18.204/30 is directly connected, Serial6/0
C       172.84.18.208/30 is directly connected, Serial3/0
R       172.84.18.212/30 [120/1] via 172.84.18.206, 00:00:04, Serial6/0
R       172.84.18.216/30 [120/2] via 172.84.18.206, 00:00:04, Serial6/0
                         [120/2] via 172.84.18.210, 00:00:22, Serial3/0
                         [120/2] via 172.84.18.193, 00:00:05, Serial2/0
R       172.84.18.220/30 [120/1] via 172.84.18.210, 00:00:22, Serial3/0
C       172.84.18.224/30 is directly connected, Serial9/0
\end{verbatim}
\end{mdframed}

    $\hfill\square$

    \vspace{12px}

    \item \textit{Realice un} traceroute \textit{del H A5 al H B6. Explica y justifica el camino que se sigue. Inidque cómo es posible que el R A3, que utiliza un protocolo de enrutamiento intra-dominio, puede obtener información de otro SA distinto que utiliza un protocolo de enrutamiento intra-dominio diferente.}
    \vspace{7px}

    \textsc{Contestación}: 
    Se puede ver que el camino que se sigue es:
    $171.84.17.1 (H A5) \xrightarrow[Fa0/1]{} 172.84.17.254 (R A5) \xrightarrow[Se2/0]{} 172.84.18.217 (R A4) \xrightarrow[Se3/0]{} 172.84.18.201 (R A1) \xrightarrow[]{} 172.84.18.205 (R A2)\xrightarrow[]{} 172.84.18.226 (R B2) \xrightarrow[]{} 171.84.17.134 (R B3) \xrightarrow[]{} 171.84.17.86 (R B6) \xrightarrow[]{} 171.84.17.1 (H B6)$

    R A3 tiene información de la organización B porque el router R B1 ha redistribuido sus rutas OSPF a RIP a R A2, asignándoles a estas una métrica de 1 en el contexto de RIP, es decir, el número de saltos
    
\begin{mdframed}
\begin{verbatim}
En H A5:

C:\>tracert 171.84.17.1

Tracing route to 171.84.17.1 over a maximum of 30 hops: 

  1   0 ms      0 ms      0 ms      172.84.17.254
  2   0 ms      10 ms     22 ms     172.84.18.217
  3   12 ms     29 ms     42 ms     172.84.18.201
  4   11 ms     72 ms     2 ms      172.84.18.205
  5   93 ms     33 ms     0 ms      172.84.18.226
  6   63 ms     72 ms     83 ms     171.84.17.134
  7   100 ms    6 ms      6 ms      171.84.17.86
  8   *         47 ms     1 ms      171.84.17.1

Trace complete.
\end{verbatim}
\end{mdframed}

    $\hfill\square$

    \vspace{12px}

    \item \textit{Tras la redistribución consulte las tablas de rutas de los routers del Área 2 para demostrar que se trata de una totally stub área. ¿Qué sucede con la tablas de rutas? ¿Por qué?}
    \vspace{7px}

    \textsc{Contestación}: Podemos observar que la tabla de rutas de R B4 no ha sufrido cambios, ya que al tratarse de un área totally stub, para llegar a cualquier IP externa a su área, debe pasar por su ABR, que es el R B3, por lo que se agrupan todas en una por defecto, incluyendo las que van a la organización A.

\begin{mdframed}
\begin{verbatim}
En R B4

Router>en
Router#sh ip ro
Router#sh ip route 
Codes: C - connected, S - static, I - IGRP, R - RIP, M - mobile, B - BGP
       D - EIGRP, EX - EIGRP external, O - OSPF, IA - OSPF inter area
       N1 - OSPF NSSA external type 1, N2 - OSPF NSSA external type 2
       E1 - OSPF external type 1, E2 - OSPF external type 2, E - EGP
       i - IS-IS, L1 - IS-IS level-1, L2 - IS-IS level-2, ia - IS-IS inter area
       * - candidate default, U - per-user static route, o - ODR
       P - periodic downloaded static route

Gateway of last resort is 171.84.16.190 to network 0.0.0.0

     171.84.0.0/26 is subnetted, 2 subnets
C       171.84.16.128 is directly connected, FastEthernet1/0
C       171.84.16.192 is directly connected, FastEthernet0/0
O*IA 0.0.0.0/0 [110/2] via 171.84.16.190, 00:05:53, FastEthernet1/0
\end{verbatim}
\end{mdframed}
    
\begin{mdframed}
\begin{verbatim}
En R B3:

Router#sh ip route
Codes: C - connected, S - static, I - IGRP, R - RIP, M - mobile, B - BGP
       D - EIGRP, EX - EIGRP external, O - OSPF, IA - OSPF inter area
       N1 - OSPF NSSA external type 1, N2 - OSPF NSSA external type 2
       E1 - OSPF external type 1, E2 - OSPF external type 2, E - EGP
       i - IS-IS, L1 - IS-IS level-1, L2 - IS-IS level-2, ia - IS-IS inter area
       * - candidate default, U - per-user static route, o - ODR
       P - periodic downloaded static route

Gateway of last resort is not set

     171.84.0.0/16 is variably subnetted, 11 subnets, 4 masks
O IA    171.84.16.0/25 [110/65] via 171.84.17.137, 00:07:11, Serial3/0
C       171.84.16.128/26 is directly connected, FastEthernet0/0
O       171.84.16.192/26 [110/2] via 171.84.16.189, 00:06:46, FastEthernet0/0
O       171.84.17.0/26 [110/65] via 171.84.17.86, 00:07:16, Serial7/0
O       171.84.17.64/28 [110/65] via 171.84.17.82, 00:07:16, Serial6/0
C       171.84.17.80/30 is directly connected, Serial6/0
C       171.84.17.84/30 is directly connected, Serial7/0
O       171.84.17.88/30 [110/95] via 171.84.17.82, 00:07:16, Serial6/0
O       171.84.17.128/30 [110/128] via 171.84.17.137, 00:07:11, Serial3/0
                         [110/128] via 171.84.17.133, 00:07:11, Serial2/0
C       171.84.17.132/30 is directly connected, Serial2/0
C       171.84.17.136/30 is directly connected, Serial3/0
     172.84.0.0/16 is variably subnetted, 12 subnets, 4 masks
O E2    172.84.16.0/23 [110/200] via 171.84.17.133, 00:07:11, Serial2/0
O E2    172.84.18.0/25 [110/200] via 171.84.17.133, 00:07:11, Serial2/0
O E2    172.84.18.128/26 [110/200] via 171.84.17.133, 00:07:11, Serial2/0
O E2    172.84.18.192/30 [110/200] via 171.84.17.133, 00:07:11, Serial2/0
O E2    172.84.18.196/30 [110/200] via 171.84.17.133, 00:07:11, Serial2/0
O E2    172.84.18.200/30 [110/200] via 171.84.17.133, 00:07:11, Serial2/0
O E2    172.84.18.204/30 [110/200] via 171.84.17.133, 00:07:11, Serial2/0
O E2    172.84.18.208/30 [110/200] via 171.84.17.133, 00:07:11, Serial2/0
O E2    172.84.18.212/30 [110/200] via 171.84.17.133, 00:07:11, Serial2/0
O E2    172.84.18.216/30 [110/200] via 171.84.17.133, 00:07:11, Serial2/0
O E2    172.84.18.220/30 [110/200] via 171.84.17.133, 00:07:11, Serial2/0
O E2    172.84.18.224/30 [110/200] via 171.84.17.133, 00:07:11, Serial2/0    
\end{verbatim}
\end{mdframed}

    $\hfill\square$

    \vspace{12px}

    \item \textit{Consulte también las tablas de rutas de los routers del Área 3 y explique por qué se trata de un área stub. ¿Qué ocurriría en el caso de que no fuera stub? ¿Por qué?}
    \vspace{7px}

    \textsc{Contestación}: Se trata de un área stub porque todos los paquetes que vayan a una red externa, la de la organización A, debe pasar por su ABR, el R B3, por lo que se agrupan estas en la ruta por defecto, pero las rutas que van a otras áreas sí aparecen, luego no es totally stub.\\
    Si no fuese stub, aparecería una entrada para cada subred, porque sí le llega información de cómo llegar a cada una de ellas, por lo que se necesita estas entradas, pero no las puede agrupar en una porque no se le indica que debe pasar siempre por su ABR, aunque este sea el caso.

\begin{mdframed}
\begin{verbatim}
En R B5:

Router>en
Router#sh ip route
Codes: C - connected, S - static, I - IGRP, R - RIP, M - mobile, B - BGP
       D - EIGRP, EX - EIGRP external, O - OSPF, IA - OSPF inter area
       N1 - OSPF NSSA external type 1, N2 - OSPF NSSA external type 2
       E1 - OSPF external type 1, E2 - OSPF external type 2, E - EGP
       i - IS-IS, L1 - IS-IS level-1, L2 - IS-IS level-2, ia - IS-IS inter area
       * - candidate default, U - per-user static route, o - ODR
       P - periodic downloaded static route

Gateway of last resort is 171.84.17.90 to network 0.0.0.0

     171.84.0.0/16 is variably subnetted, 10 subnets, 4 masks
O IA    171.84.16.0/25 [110/127] via 171.84.17.90, 00:14:11, Serial3/0
O IA    171.84.16.128/25 [110/64] via 171.84.17.90, 00:13:21, Serial3/0
O       171.84.17.0/26 [110/32] via 171.84.17.90, 00:14:11, Serial3/0
C       171.84.17.64/28 is directly connected, FastEthernet0/0
C       171.84.17.80/30 is directly connected, Serial2/0
O       171.84.17.84/30 [110/62] via 171.84.17.90, 00:14:11, Serial3/0
C       171.84.17.88/30 is directly connected, Serial3/0
O IA    171.84.17.128/30 [110/190] via 171.84.17.90, 00:14:11, Serial3/0
O IA    171.84.17.132/30 [110/126] via 171.84.17.90, 00:14:11, Serial3/0
O IA    171.84.17.136/30 [110/126] via 171.84.17.90, 00:14:11, Serial3/0
O*IA 0.0.0.0/0 [110/63] via 171.84.17.90, 00:14:11, Serial3/0
\end{verbatim}
\end{mdframed}


\begin{mdframed}
\begin{verbatim}
En R B6:

Router>en
Router#sh ip route
Codes: C - connected, S - static, I - IGRP, R - RIP, M - mobile, B - BGP
       D - EIGRP, EX - EIGRP external, O - OSPF, IA - OSPF inter area
       N1 - OSPF NSSA external type 1, N2 - OSPF NSSA external type 2
       E1 - OSPF external type 1, E2 - OSPF external type 2, E - EGP
       i - IS-IS, L1 - IS-IS level-1, L2 - IS-IS level-2, ia - IS-IS inter area
       * - candidate default, U - per-user static route, o - ODR
       P - periodic downloaded static route

Gateway of last resort is 171.84.17.85 to network 0.0.0.0

     171.84.0.0/16 is variably subnetted, 10 subnets, 4 masks
O IA    171.84.16.0/25 [110/96] via 171.84.17.85, 00:15:00, Serial2/0
O IA    171.84.16.128/25 [110/33] via 171.84.17.85, 00:14:10, Serial2/0
C       171.84.17.0/26 is directly connected, FastEthernet0/0
O       171.84.17.64/28 [110/65] via 171.84.17.89, 00:15:10, Serial3/0
O       171.84.17.80/30 [110/95] via 171.84.17.85, 00:15:20, Serial2/0
C       171.84.17.84/30 is directly connected, Serial2/0
C       171.84.17.88/30 is directly connected, Serial3/0
O IA    171.84.17.128/30 [110/159] via 171.84.17.85, 00:15:00, Serial2/0
O IA    171.84.17.132/30 [110/95] via 171.84.17.85, 00:15:10, Serial2/0
O IA    171.84.17.136/30 [110/95] via 171.84.17.85, 00:15:10, Serial2/0
O*IA 0.0.0.0/0 [110/32] via 171.84.17.85, 00:15:10, Serial2/0
\end{verbatim}
\end{mdframed}

    $\hfill\square$
\end{enumerate}

\end{document}
